\documentclass[12pt]{article}
usepackage{amsmath}
\usepackage{amssymb}
\usepackage{amsmath,amssymb,amsthm}
\usepackage{accessibility}
\usepackage{geometry}
\geometry{margin=1in}

\newtheorem{theorem}{Theorem}
\newtheorem{lemma}[theorem]{Lemma}
\newtheorem{definition}[theorem]{Definition}
\newtheorem{proposition}[theorem]{Proposition}

\title{The Unified Algebraic-Lattice Bipartition Framework\\
\large A Synthesis of ALCF and AMBS for Quasiperfect Number Search}
\author{math\_explorer Project}
\date{\today}

\begin{document}
\maketitle

\begin{abstract}
We present the Unified Algebraic-Lattice Bipartition Framework (UALBF), a
four-phase computational sieve that combines the Algebraic-Lattice Convergence
Framework (ALCF) and the Algebraic-Modular Bipartition Sieve (AMBS) to search
for quasiperfect numbers. A quasiperfect number $N$ satisfies $\sigma(N) = 2N+1$,
where $\sigma$ denotes the sum-of-divisors function. No quasiperfect number is
currently known; UALBF provides a structured pipeline to establish lower bounds
and explore the algebraic constraints governing their possible existence.
\end{abstract}

\section{Introduction}

A \emph{perfect number} $N$ satisfies $\sigma(N) = 2N$. Relaxing this identity
by one yields a \emph{quasiperfect number}: an integer $N$ with $\sigma(N) = 2N+1$.
Despite extensive computational searches, no quasiperfect number has ever been
found \cite{Guy2004}. It is known that any quasiperfect number must be an odd
perfect square exceeding $10^{35}$ \cite{CohenHagis1982}.

UALBF is motivated by the observation that both the ALCF (which studies the
multiplicative structure of $\sigma$) and the AMBS (which applies modular
arithmetic sieves to factor products) provide complementary constraints on any
hypothetical quasiperfect number. By unifying these two frameworks, UALBF
achieves superior pruning compared to either approach alone.

\section{Mathematical Preliminaries}

\begin{definition}[Sum of Divisors]
For $n \geq 1$, the \emph{sum-of-divisors function} is
\[
    \sigma(n) = \sum_{d \mid n} d.
\]
For prime powers, $\sigma(p^k) = 1 + p + p^2 + \cdots + p^k = \frac{p^{k+1}-1}{p-1}$.
\end{definition}

\begin{definition}[Quasiperfect Number]
A positive integer $N$ is \emph{quasiperfect} if $\sigma(N) = 2N + 1$.
\end{definition}

Since $\sigma$ is multiplicative and $N$ must be an odd square (see \cite{CohenHagis1982}),
we write $N = \prod_i p_i^{2e_i}$ and note that
\[
    \sigma(N) = \prod_i \sigma(p_i^{2e_i}).
\]

\section{The Four-Phase Pipeline}

\subsection{Phase 1: Global Annihilation Sieve}

For each odd prime $p \leq L$ and exponent $1 \leq e \leq E$, we compute
$\sigma(p^{2e})$ and factor it. A component $p^{2e}$ is \emph{annihilated} if
$\sigma(p^{2e})$ has a prime divisor $q$ with $q \equiv 5 \pmod{8}$ or
$q \equiv 7 \pmod{8}$.

\begin{proposition}
If $q \mid \sigma(p^{2e})$ and $q \equiv 5$ or $7 \pmod{8}$, then $p^{2e}$
cannot appear in the factorization of a quasiperfect number.
\end{proposition}

\begin{proof}
By quadratic reciprocity and properties of the Jacobi symbol, such a prime
$q$ forces $2N+1 \not\equiv 0 \pmod{q}$, contradicting $\sigma(N) = 2N+1$
when $q \mid \sigma(N)$. See \cite{CohenHagis1982} for the full argument.
\end{proof}

\subsection{Phase 2: Prefix Tree Construction}

Let $\mathcal{C}$ be the set of components surviving Phase~1. A \emph{prefix}
is a pair $(N_L, S_L)$ where
\[
    N_L = \prod_{i \in I} p_i^{2e_i}, \qquad
    S_L = \prod_{i \in I} \sigma(p_i^{2e_i}),
\]
for some index set $I \subseteq \mathcal{C}$. We build all prefixes with
$N_L \geq T$ (a chosen threshold) by a depth-first expansion of the
factor lattice.

\subsection{Phase 3: Lattice Oracle}

For each prefix $(N_L, S_L)$, a necessary condition for extension to a
quasiperfect number is the existence of an integer $x_L$ satisfying
\[
    -2 N_L \cdot x_L \equiv 1 \pmod{S_L}.
\]
If $\gcd(2N_L, S_L) \nmid 1$, no such $x_L$ exists and the prefix is
\emph{rejected}.

\subsection{Phase 4: Exact Ray Casting}

For surviving prefixes, we write $N = N_L \cdot z^2$ and require
$S_L \cdot \sigma(z^2) = 2 N_L z^2 + 1$. Setting $x_L \equiv (-2N_L)^{-1} \pmod{S_L}$,
we compute modular square roots of $x_L$ modulo $S_L$ via the AMBS and enumerate
candidate values of $z$, checking whether $\sigma(N) = 2N+1$ exactly.

\section{Implementation}

The UALBF is implemented in Rust as part of the \texttt{math\_explorer} crate
under \texttt{pure\_math::number\_theory::ualbf}. The pipeline is fully
parameterized by:
\begin{itemize}
    \item \texttt{limit\_p}: upper bound on primes in Phase~1,
    \item \texttt{max\_e}: maximum exponent in Phase~1,
    \item \texttt{stop\_threshold}: minimum prefix product $N_L$,
    \item \texttt{target\_max}: upper bound on candidates in Phase~4.
\end{itemize}

\section{Complexity and Expected Behaviour}

The number of surviving components after Phase~1 is typically a small fraction
of all prime powers up to $L$. Empirically, primes $p \equiv 3 \pmod{8}$ tend
to survive because $\sigma(p^2) = p^2 + p + 1$ often avoids the annihilating
residue classes. Phase~3 provides additional pruning through modular
invertibility constraints, and Phase~4 terminates quickly in the absence of
quasiperfect numbers.

\section{Conclusion}

UALBF provides a unified, modular, and extensible framework for the
computational study of quasiperfect numbers. Future work includes
parallelising the prefix enumeration, extending the sieve with additional
congruence conditions, and integrating with certified arithmetic libraries
for rigorous lower-bound proofs.

\bibliographystyle{plain}
\bibliography{references}

\end{document}
